% ===================================================================
%  Αρχείο: 21854.tex (Fragment)
%  Αυτόματη μετατροπή μέσω doc2latex.py
% ===================================================================

ΘΕΜΑ 4

Ένα δοχείο περιέχει υγρό το οποίο εξατμίζεται. Στο παρακάτω διάγραμμα φαίνεται η ποσότητα $Q$, σε λίτρα, του υγρού που έχει απομείνει στο δοχείο μετά από $t$ ημέρες.


% TODO: Μετατροπή σχήματος σε TikZ/PGFPlots
\begin{center}

\begin{tikzpicture}[scale=0.8]
    % Grid
    \draw[help lines, gray!30] (0,0) grid (8,10);
    
    % Axes
    \draw[->, thick] (-0.5,0) -- (8.5,0) node[right] {$t \text{ ημέρες}$};
    \draw[->, thick] (0,-0.5) -- (0,10.5) node[above] {$Q(t) \text{ (lt)}$};
    
    % Function Q(t) = 8 * 2^{-t/2}
    \draw[ultra thick, blue, domain=0:7.5, samples=100] plot (\x, {8*2^(-\x/2)});
    
    % Key Points
    \filldraw (0,8) circle (2.5pt) node[left] {$8$};
    \foreach \t/\q in {2/4, 4/2} {
        \filldraw (\t,\q) circle (2pt);
        \draw[dashed] (\t,\q) -- (\t,0) node[below] {$\t$};
        \draw[dashed] (\t,\q) -- (0,\q) node[left] {$\q$};
    }
    
    % Origin
    \node[below left] at (0,0) {$0$};

\end{tikzpicture}

% % \caption{Σχήμα 1}
\end{center}


Η ποσότητα του υγρού στο δοχείο μειώνεται εκθετικά και μετά από $t$ ημέρες δίνεται από τη σχέση $Q(t) = Q_{0}2^{- \frac{t}{c}},\ c\mathbb{\in R}$, όπου $Q_{0}$ η αρχική ποσότητα του υγρού.

\begin{alist}
\item Με βάση το διάγραμμα:

\end{alist}
\begin{alist}
\item
  να συμπληρώσετε τον παρακάτω πίνακα:

\begin{longtable}[]{@{}
  >{\raggedright\arraybackslash}p{(\linewidth - 8\tabcolsep) * \real{0.5980}}
  >{\raggedright\arraybackslash}p{(\linewidth - 8\tabcolsep) * \real{0.0894}}
  >{\raggedright\arraybackslash}p{(\linewidth - 8\tabcolsep) * \real{0.1013}}
  >{\raggedright\arraybackslash}p{(\linewidth - 8\tabcolsep) * \real{0.0996}}
  >{\raggedright\arraybackslash}p{(\linewidth - 8\tabcolsep) * \real{0.1117}}@{}}
\noalign{}
\begin{minipage}[b]{\linewidth}\raggedright
Χρόνος $t$ σε ημέρες
\end{minipage} & \begin{minipage}[b]{\linewidth}\raggedright
\begin{equation*}0\end{equation*}
\end{minipage} & \begin{minipage}[b]{\linewidth}\raggedright
\begin{equation*}2\end{equation*}
\end{minipage} & \begin{minipage}[b]{\linewidth}\raggedright
\begin{equation*}4\end{equation*}
\end{minipage} & \begin{minipage}[b]{\linewidth}\raggedright
\begin{equation*}6\end{equation*}
\end{minipage} \\
\noalign{}

\noalign{}
\endlastfoot
Ποσότητα $Q(t)$ του υγρού σε λίτρα. & & & & \\
\end{longtable}

\monades{5}


\end{alist}
\begin{alist}
\item
  να βρείτε την αρχική ποσότητα $Q_{0}$ του υγρού,

\monades{3}


\end{alist}
\begin{alist}
\item
  να βρείτε το χρόνο που χρειάζεται για να εξατμιστεί η μισή ποσότητα του υγρού που υπήρχε τη χρονική στιγμή $t = 0$ στο δοχείο.

\monades{5}

\item Αν $Q_{0} = 8$ και $Q(2) = 4$, να δείξετε ότι $c = 2$.

\monades{7}

\item Αν $Q(t) = 8 \cdot 2^{- \frac{t}{2}}$, να δείξετε ότι χρειάζεται να περάσουν δύο ημέρες για να εξατμιστεί η μισή ποσότητα $Q(t)$ του υγρού που υπάρχει στο δοχείο οποιαδήποτε χρονική στιγμή $t$.

\monades{5}
\end{alist}
